% GENERATED FILE. Edit canonical lessons, then run npm run generate:books.
% canonical-source-hash: fnv1a64:c7cdadbaaf2e131f
% canonical-lessons: AR-C500-maghrur, AR-C500-mutakabbir, AR-C500-anani, AR-C500-jaban, AR-C500-jari

\chapter{Vain, selfish, cowardly and bold}
\label{ch:ar-vain-selfish-cowardly-and-bold}
\hlchaptermodality{500}

\begin{chapteropening}
\textbf{By the end of this chapter:} \emph{I can describe a person as conceited, arrogant, selfish, cowardly or bold.}

Complete the last lesson of chapter 500: I can describe a person as conceited, arrogant, selfish, cowardly or bold.
\end{chapteropening}

\section[\texorpdfstring{\ar{مغرور} (maghrūr)}{مغرور (maghrūr)}]{\ar{مغرور} (maghrūr) — conceited, vain}

\label{lesson:AR-C500-maghrur}



\begin{quote}
Before the new one: say the Arabic for loyal, then the Arabic for faithful.
\end{quote}

\subsection*{You'll want to know: \ar{مغرور}}
\textbf{\ar{مغرور}} — \emph{maghrūr} — \textquotedblleft{}conceited, vain\textquotedblright{}.

\begin{quote}
\ar{لا} \ar{يحب} \ar{الفريق} \ar{اللاعب} \ar{المغرور}.
\end{quote}

\begin{quote}
\emph{lā yuḥibbu al-farīq al-lāʿib al-maghrūr.}
\end{quote}

\begin{quote}
The team does not like the conceited player.
\end{quote}

\begin{cousinweb}
\textbf{\ar{مغرور}} (\emph{maghrūr}) is built on the native Arabic root \emph{gh-r-r}, \textquotedblleft{}to deceive, to beguile\textquotedblright{} (Lane). Maghrūr is the passive participle: one beguiled, here by his own worth.
\end{cousinweb}

\begin{grammarlens}[title={What you've built}]
One more word for describing what you lived and felt.
\end{grammarlens}

\subsection*{Your turn}
\begin{itemize}
\raggedright
  \item \emph{Say it:} \emph{maghrūr}
  \item \emph{Say it:} \emph{maghrūr}, once more
  \item \emph{Say it:} \emph{lā yuḥibbu al-farīq al-lāʿib al-maghrūr.} — The team does not like the conceited player.
  \item \emph{Recall:} say the Arabic for loyal, then the Arabic for faithful, then say \emph{maghrūr} again
\end{itemize}

\begin{tcolorbox}[breakable,colback=teal!4,colframe=teal!35!black,title={Before you move on}]
What does \ar{مغرور} mean? (\textquotedblleft{}Conceited, vain\textquotedblright{}.) Say it once more.
\end{tcolorbox}
\section[\texorpdfstring{\ar{متكبر} (mutakabbir)}{متكبر (mutakabbir)}]{\ar{متكبر} (mutakabbir) — arrogant}

\label{lesson:AR-C500-mutakabbir}



\begin{quote}
Before the new one: say the Arabic for faithful, then the Arabic for conceited.
\end{quote}

\subsection*{You'll want to know: \ar{متكبر}}
\textbf{\ar{متكبر}} — \emph{mutakabbir} — \textquotedblleft{}arrogant\textquotedblright{}.

\begin{quote}
\ar{كان} \ar{المدير} \ar{متكبرا} \ar{مع} \ar{الموظفين}.
\end{quote}

\begin{quote}
\emph{kāna al-mudīr mutakabbiran maʿa al-muwaẓẓafīn.}
\end{quote}

\begin{quote}
The manager was arrogant with the staff.
\end{quote}

\begin{cousinweb}
\textbf{\ar{متكبر}} (\emph{mutakabbir}) is built on the native Arabic root \emph{k-b-r}, \textquotedblleft{}to be big, great\textquotedblright{} (Lane). Kabīr, big, is from the same root; Form V mutakabbir is one who makes himself big.
\end{cousinweb}

\begin{grammarlens}[title={What you've built}]
One more word for describing what you lived and felt.
\end{grammarlens}

\subsection*{Your turn}
\begin{itemize}
\raggedright
  \item \emph{Say it:} \emph{mutakabbir}
  \item \emph{Say it:} \emph{mutakabbir}, once more
  \item \emph{Say it:} \emph{kāna al-mudīr mutakabbiran maʿa al-muwaẓẓafīn.} — The manager was arrogant with the staff.
  \item \emph{Recall:} say the Arabic for faithful, then the Arabic for conceited, then say \emph{mutakabbir} again
\end{itemize}

\begin{tcolorbox}[breakable,colback=teal!4,colframe=teal!35!black,title={Before you move on}]
What does \ar{متكبر} mean? (\textquotedblleft{}Arrogant\textquotedblright{}.) Say it once more.
\end{tcolorbox}
\section[\texorpdfstring{\ar{أناني} (anānī)}{أناني (anānī)}]{\ar{أناني} (anānī) — selfish}

\label{lesson:AR-C500-anani}



\begin{quote}
Before the new one: say the Arabic for conceited, then the Arabic for arrogant.
\end{quote}

\subsection*{You'll want to know: \ar{أناني}}
\textbf{\ar{أناني}} — \emph{anānī} — \textquotedblleft{}selfish\textquotedblright{}.

\begin{quote}
\ar{الطفل} \ar{الأناني} \ar{لا} \ar{يشارك} \ar{أخاه} \ar{في} \ar{ألعابه}.
\end{quote}

\begin{quote}
\emph{aṭ-ṭifl al-anānī lā yushāriku akhāhu fī alʿābihi.}
\end{quote}

\begin{quote}
The selfish child does not share his toys with his brother.
\end{quote}

\begin{cousinweb}
\textbf{\ar{أناني}} (\emph{anānī}) is a modern Arabic word, usually traced to the pronoun \emph{anā}, \textquotedblleft{}I\textquotedblright{}, with the adjective ending \emph{-ī}: a person always saying I. The noun \emph{anāniyya} is selfishness.
\end{cousinweb}

\begin{grammarlens}[title={What you've built}]
One more word for describing what you lived and felt.
\end{grammarlens}

\subsection*{Your turn}
\begin{itemize}
\raggedright
  \item \emph{Say it:} \emph{anānī}
  \item \emph{Say it:} \emph{anānī}, once more
  \item \emph{Say it:} \emph{aṭ-ṭifl al-anānī lā yushāriku akhāhu fī alʿābihi.} — The selfish child does not share his toys with his brother.
  \item \emph{Recall:} say the Arabic for conceited, then the Arabic for arrogant, then say \emph{anānī} again
\end{itemize}

\begin{tcolorbox}[breakable,colback=teal!4,colframe=teal!35!black,title={Before you move on}]
What does \ar{أناني} mean? (\textquotedblleft{}Selfish\textquotedblright{}.) Say it once more.
\end{tcolorbox}
\section[\texorpdfstring{\ar{جبان} (jabān)}{جبان (jabān)}]{\ar{جبان} (jabān) — cowardly; a coward}

\label{lesson:AR-C500-jaban}



\begin{quote}
Before the new one: say the Arabic for arrogant, then the Arabic for selfish.
\end{quote}

\subsection*{You'll want to know: \ar{جبان}}
\textbf{\ar{جبان}} — \emph{jabān} — \textquotedblleft{}cowardly; a coward\textquotedblright{}.

\begin{quote}
\ar{هرب} \ar{الرجل} \ar{الجبان} \ar{من} \ar{الكلب}.
\end{quote}

\begin{quote}
\emph{haraba ar-rajul al-jabān min al-kalb.}
\end{quote}

\begin{quote}
The cowardly man ran away from the dog.
\end{quote}

The noun \ar{جبن}, cowardice, is written like jubn, cheese.

\begin{cousinweb}
\textbf{\ar{جبان}} (\emph{jabān}) is built on the native Arabic root \emph{j-b-n}, \textquotedblleft{}to be cowardly\textquotedblright{} (Lane). Jabān is the adjective of a lasting trait.
\end{cousinweb}

\begin{grammarlens}[title={What you've built}]
One more word for describing what you lived and felt.
\end{grammarlens}

\subsection*{Your turn}
\begin{itemize}
\raggedright
  \item \emph{Say it:} \emph{jabān}
  \item \emph{Say it:} \emph{jabān}, once more
  \item \emph{Say it:} \emph{haraba ar-rajul al-jabān min al-kalb.} — The cowardly man ran away from the dog.
  \item \emph{Recall:} say the Arabic for arrogant, then the Arabic for selfish, then say \emph{jabān} again
\end{itemize}

\begin{tcolorbox}[breakable,colback=teal!4,colframe=teal!35!black,title={Before you move on}]
What does \ar{جبان} mean? (\textquotedblleft{}Cowardly; a coward\textquotedblright{}.) Say it once more.
\end{tcolorbox}
\section[\texorpdfstring{\ar{جريء} (jarīʾ)}{جريء (jarīʾ)}]{\ar{جريء} (jarīʾ) — bold, daring}

\label{lesson:AR-C500-jari}



\begin{quote}
Before the new one: say the Arabic for selfish, then the Arabic for cowardly.
\end{quote}

\subsection*{You'll want to know: \ar{جريء}}
\textbf{\ar{جريء}} — \emph{jarīʾ} — \textquotedblleft{}bold, daring\textquotedblright{}.

\begin{quote}
\ar{قفز} \ar{الولد} \ar{الجريء} \ar{في} \ar{النهر}.
\end{quote}

\begin{quote}
\emph{qafaza al-walad al-jarīʾ fī an-nahr.}
\end{quote}

\begin{quote}
The bold boy jumped into the river.
\end{quote}

\begin{cousinweb}
\textbf{\ar{جريء}} (\emph{jarīʾ}) is built on the native Arabic root \emph{j-r-ʾ}, \textquotedblleft{}to be bold, to dare\textquotedblright{} (Lane). Jarīʾ is the adjective pattern of a lasting trait.
\end{cousinweb}

\begin{grammarlens}[title={What you've built}]
One more word for describing what you lived and felt.
\end{grammarlens}

\subsection*{Your turn}
\begin{itemize}
\raggedright
  \item \emph{Say it:} \emph{jarīʾ}
  \item \emph{Say it:} \emph{jarīʾ}, once more
  \item \emph{Say it:} \emph{qafaza al-walad al-jarīʾ fī an-nahr.} — The bold boy jumped into the river.
  \item \emph{Recall:} say the Arabic for selfish, then the Arabic for cowardly, then say \emph{jarīʾ} again
\end{itemize}

\begin{tcolorbox}[breakable,colback=teal!4,colframe=teal!35!black,title={Before you move on}]
What does \ar{جريء} mean? (\textquotedblleft{}Bold, daring\textquotedblright{}.) Say it once more.
\end{tcolorbox}
